\documentclass[exercice,a4]{thedocument}%
\usepackage{texmaths-tikz}%
\usepackage{pstricks}%
\usepackage{pst-bezier}%
\usepackage{pstricks-add}%
\usepackage{pst-eucl}%
\usepackage{pst-func}%
\usepackage{pst-tree}%
\usepackage{pst-circ}%
\usepackage{tikz}%
\usepackage{pgfplots}%
\usepackage{tkz-tab}%
\usepackage{tikzpagenodes}%
\usepackage{tabularray}%
\usepackage{tabularx}%
\usepackage{multirow}%
\usepackage{multicol}%
\usepackage{scratch3}%
\usepackage{enumitem}%
\usepackage[np]{numprint}%
\usepackage{pas-tableur}%
\usepackage{listings}%
\usepackage{graphicx}%
\usepackage{adjustbox}%
\usepackage{needspace}%
\usepackage{xcolor}%
\usepackage{eurosym}%
\usepackage{diagbox}%
\usepackage{colortbl}%
\usepackage{multido}%
\usepackage{icomma}%
\usepackage[french]{babel}%
\startdatakeys
    \source{26GENMATMEAG1}
    \BO{2026}
    \cycle{4}
    \competencesDuSocle{CA,RA,CO}
    \objectifsApprentissage{C4A1aa,C4A2ad,C4B2cb,C4D2aa,C4D3ac}
    \motsCles{pourcentage,médiane}
\enddatakeys
\definecolor{gris}{gray}{0.5}%
\providecommand{\wideparen}[1]{\overset{\frown}{#1}}%
\begin{document}%
\providecommand{\degres}{^\circ}%
Le tableau ci-dessous présente le nombre de médailles obtenues par 9 pays lors des Jeux Paralympiques de Paris 2024.

\begin{center}
	\begin{tblr}{hlines, vlines, colspec={X[l,4cm]*{4}{X[c,2cm]}}}
		Pays            & Or & Argent & Bronze & Total \\
		Chine           & 94 & 76     & 50     & 220   \\
		Grande-Bretagne & 49 & 44     & 31     & 124   \\
		États-Unis      & 36 & 42     & 27     & 105   \\
		Pays-Bas        & 27 & 17     & 12     & 56    \\
		Brésil          & 25 & 26     & 38     & 89    \\
		Italie          & 24 & 15     & 32     & 71    \\
		Ukraine         & 22 & 28     & 32     & 82    \\
		France          & 19 & 28     & 28     & 75    \\
		Australie       & ?  & 17     & 28     & 63
	\end{tblr}
\end{center}

\begin{enumerate}
	\item Combien de médailles d'or ont été obtenues par les Pays-Bas?

	\item Calculer le nombre de médailles d'or obtenues par l'Australie.

	\item Montrer que l'affirmation suivante est vraie.

	« Plus de $20~\%$ des médailles obtenues par la Grande-Bretagne sont en bronze.»

	\item \begin{enumerate}
		\item Déterminer la médiane de la série des nombres de médailles obtenues par ces 9 pays (colonne « Total » du tableau).

		\textbf{Détailler la réponse en précisant la démarche.}

		\item Interpréter ce résultat dans le contexte de l'exercice.
	\end{enumerate}

	\item Aux Jeux Paralympiques de Tokyo en 2021, le Brésil avait obtenu 20 médailles d'argent. Pour ceux de Paris en 2024, il a obtenu 26 médailles d'argent.

	Quel est le pourcentage d'augmentation du nombre de médailles d'argent obtenues par le Brésil entre 2021 et 2024 ?
\end{enumerate}

\needspace{10\baselineskip}
% saut de page si moins de 10 lignes disponibles
\medskip
\showthesource
\end{document}
